\documentclass[10pt,a4paper]{amsart}
\usepackage[margin=19mm]{geometry}
\usepackage{amsmath,amssymb,mathtools}
\usepackage[scr=rsfs,cal=euler]{mathalfa}
\usepackage{xfrac}
\usepackage[unicode,hidelinks,bookmarks=false]{hyperref}
\newcommand{\ie}{i.e.\ }
\newcommand{\cf}{cf.\ }
\newcommand{\resp}{respectively\ }
\newcommand{\Subsection}{\subsection}
\providecommand{\nobreakdash}{\nobreak\hbox{-}}
\setlength{\emergencystretch}{2em}
\multlinegap0pt
\usepackage{bm}
\usepackage{bbm}
\usepackage{dsfont}
\usepackage{xspace}
\usepackage{cancel}
\usepackage{mathtools}
\usepackage{amsthm}
\makeatletter
\@ifundefined{lemma}{\newtheorem{lemma}{Lemma}}{}
\@ifundefined{notation}{\newtheorem{notation}{Notation}}{}
\@ifundefined{theorem}{\newtheorem{theorem}{Theorem}}{}
\@ifundefined{proposition}{\newtheorem{proposition}[theorem]{Proposition}}{}
\makeatother
\newcommand{\E}{\mathbb{E}}
\newcommand{\I}{\mathds{1}}
\newcommand{\NC}{\mathcal{NC}}
\newcommand{\Tr}{\text{Tr}}
\newcommand{\cP}{\mathcal{P}}
\title[Asymptotic infinitesimal freeness of covariance matrices]{Asymptotic infinitesimal freeness of covariance matrices}
\author{Daniel Munoz George, Pei-Lun Tseng}
\date{Source: arXiv:2609.18187v1 (CC BY 4.0)}
\begin{document}
\maketitle

\noindent\emph{Source and licence.} Asymptotic infinitesimal freeness of covariance matrices. Version arXiv:2609.18187v1 (CC BY 4.0), distributed under the Creative Commons Attribution 4.0 International licence, as recorded in the arXiv licence metadata for this exact version and shown on the abstract page https://arxiv.org/abs/2609.18187v1. Full rights notice: licenses/arxiv-2609.18187-CC-BY-4.0.txt.\par\smallskip

\noindent\textit{Editorial scope.} Counted unit: the Proposition whose source label is \texttt{Proposition: Infinitesimal contribution pi=q for 2-4 tree} in the article (source0.tex lines 1737--1744, source label \texttt{Proposition: Infinitesimal contribution pi=q for 2-4 tree}), delivered together with the article's own complete original proof of it (source0.tex lines 1745--1780, one \texttt{proof} environment of 36 lines). No mathematics is written, repaired or re-derived: statement and proof are the article's own text line for line. Required context retained uncounted: Lemma: Expansion of moments 4 (lines 349--354); Lemma: Expansion of moments 5 (lines 436--444); Notation: New pi of 2-4 tree (lines 1332--1336); Lemma: Characterization of 2-4 tree (lines 1351--1353); Proposition: Equivalence of ker(r) and Ker(rr) (lines 1502--1504); Proposition: contribution of 2-4 tree (lines 1593--1599); Notation: Distance in double tree (lines 1678--1735). Every definition, display and label of those lines is retained unchanged. Omitted from this package and not redistributed: 1--348, 355--435, 445--1331, 1337--1350, 1354--1501, 1505--1592, 1600--1677, 1736--1736, 1781--3291. Nothing inside a retained excerpt is omitted or altered: every retained body is delivered exactly as the article's lines are written, so no omission receipt of any kind accompanies this package. Citations: the retained excerpts carry no citation command at all, so no bibliography entry of the article is transcribed and no reference is redistributed. Reference adaptation: none was needed and none was made; every reference inside the delivered unit resolves inside it, so no unresolved reference or citation is produced. Printed slips kept literal: no printed word, symbol, hypothesis or number of the counted statement or of its proof was altered. Layout and class: the article's own class is \texttt{article} with packages including graphicx, amsmath, amssymb, amsthm, xcolor, dsfont, hyperref, geometry, tikz, sansmath, caption, authblk; the wrapper is the lane's pinned \texttt{amsart} editorial wrapper at 10pt on A4, which restates the theorem-like environments the retained text needs and adds the article's own macro definitions that the retained text needs (\texttt{\textbackslash E}, \texttt{\textbackslash I}, \texttt{\textbackslash NC}, \texttt{\textbackslash Tr}, \texttt{\textbackslash cP}), with extra wrapper packages (bm, bbm, dsfont, xspace, cancel, mathtools). The delivered numbering is the unit's own. Assets: no retained span contains a figure environment, a graphics-inclusion command, a PDF-page inclusion, a table or a TikZ picture; no image file is copied into this package, the package declares no asset dependency and no image file is redistributed with it. Licence: version arXiv:2609.18187v1 of this article is distributed under the Creative Commons Attribution 4.0 International licence, as recorded in the arXiv licence metadata for this exact version and shown on the abstract page https://arxiv.org/abs/2609.18187v1. A full rights notice is delivered at licenses/arxiv-2609.18187-CC-BY-4.0.txt. Characters: every retained excerpt is pure ASCII, so the delivered engine, XeLaTeX with its default Latin Modern text font, reports no missing character.
\begin{theorem}\label{Lemma: Expansion of moments 4}
Let $M^{(r_1)},\dots, M^{(r_q)}$ be a collection of independent covariance matrices, then
\begin{equation*}
\frac{1}{n}\E\circ \textup{\Tr}(M^{(r_1)}\cdots M^{(r_q)})=\sum_{\substack{\sigma\in \NC(q) \\ \sigma\leq ker(r)}}\left(\frac{p}{n}\right)^{\#(\sigma)}+O(\frac{1}{n}).
\end{equation*}
\end{theorem}

\begin{theorem}\label{Lemma: Expansion of moments 5}
Let $M^{(r_1)},\dots, M^{(r_q)}$ be a collection of independent covariance matrices such that $\E(x_{1,2}^2)=0$, then
\begin{align*}
&\frac{1}{n}\E\circ \textup{\Tr}(M^{(r_1)}\cdots M^{(r_q)})\\
&=\sum_{\substack{\sigma\in \NC(q) \\ \sigma\leq ker(r)}}\left(\frac{p}{n}\right)^{\#(\sigma)}+ 
\frac{1}{n}\sum_{\substack{\sigma\in \NC(q) \\ \sigma\leq ker(r)}}\left(\frac{p}{n}\right)^{\#(\sigma)-1}\Bigg(\sum_{\substack{A,B\in 2\sigma \\ d_{T^\rho}(A,B)=2}}(\alpha_{a,b}-1)+\sum_{\substack{A,B\in 2Kr(\sigma)-1 \\ d_{T^\rho}(A,B)=2}}\frac{p}{n}(\alpha_{a,b}-1)\Bigg) +O(\frac{1}{n^2}).
\end{align*}
Here $\alpha_{a,b}=\frac{1}{2}\E|x_{1,2}|^4 \I_{r_{\frac{a}{2}}=r_{\frac{b}{2}}}+\I_{r_{\frac{a}{2}}\neq r_{\frac{b}{2}}}$, where $a,b$ are even numbers that depend on $A,B$ and are defined as in Notation \ref{Notation: Distance in double tree}. Further, $2\sigma$ is the partition whose blocks are obtained from multiplying each element of the blocks of $\sigma$ by 2 while $2Kr(\sigma)-1$ is the partition whose blocks are obtained from multiplying each element of the blocks of $Kr(\sigma)$ by 2 and subtracting 1. Finally $\rho=2\sigma\cup (2Kr(\sigma)-1)$.
\end{theorem}

\begin{notation}\label{Notation: New pi of 2-4 tree}
Let $\pi\in \cP(m)$ be such that $T^\pi$ is a 2-4 tree. Let 
$$B=\{e_u,e_v,e_w,e_x\}$$ 
be the edge of $\overline{T^\pi}$ of multiplicity $4$ with $e_u,e_w$ having the same orientation and opposite to $e_v,e_x$. Let $\overline{\pi}^\prime\in \cP_2(m)$ be the pairing whose blocks are the blocks of $\overline{\pi}$ except for $B$ which turns into the blocks $\{u,v\}$ and $\{w,x\}$ of $\overline{\pi}^\prime$.
\end{notation}

\begin{lemma}\label{Lemma: Characterization of 2-4 tree}
Let $\pi\in \cP(m)$. Then $T^\pi$ is a 2-4 tree if and only if $T^\pi$ is obtained by identifying two vertices of $T^\sigma$ that are at distance $2$, where $T^\sigma$ is a double tree.
\end{lemma}

\begin{proposition}\label{Proposition: Equivalence of ker(r) and Ker(rr)}
Let $\sigma\in \NC(\{2,4,\dots,2q\})$ and $Kr(\sigma)\in \NC(\{1,3,\dots,2q-1\})$ be its kreweras complement. Let $\pi=\sigma\cup Kr(\sigma)$ and let $\overline{\pi}$ be such that $\pi=\gamma\overline{\pi}$, with $\gamma=(1,\dots,2q)$. Let $r_1,\dots,r_q$ be a choice of indices. Then $\overline{\pi}\leq ker(rr)$ if and only if $\sigma/2 \leq ker(r)$. Here $\sigma/2$ is the partition of $\NC(q)$ obtained by dividing every element by two of the blocks of $\sigma$.
\end{proposition}

\begin{proposition}\label{Proposition: contribution of 2-4 tree}
Let $M^{(r_1)},\dots, M^{(r_q)}$ be a collection of independent covariance matrices. Let $\pi\in \cP(2q)$ be such that $C^\pi$ is a 2-4 tree. Let $i_1,\dots,i_{2q}$ be such that $ker(i)=\pi$. If $\E(x_i)\neq 0$, then either
\begin{enumerate}
    \item $\overline{\pi}^\prime \leq ker(rr)$ and $\E(x_i)=1$, further only one of the choices for $\overline{\pi}^\prime$ satisfies $\overline{\pi}^\prime \leq ker(rr)$. Here $\overline{\pi}^\prime$ is as in Notation \ref{Notation: New pi of 2-4 tree}, or,
    \item $\overline{\pi}\leq ker(rr)$ and $\E(x_i)=\E|x_{1,2}|^4$. 
\end{enumerate}
\end{proposition}

\begin{notation}\label{Notation: Distance in double tree}
Let $\rho\in \cP(2q)$ be such that $T^\rho$ is a double tree. 
\begin{enumerate}
    \item Given any two blocks $A,B\in \rho$ that correspond to vertices of $\overline{T^\rho}$, there is a unique path from $A$ to $B$ in $\overline{T^\rho}$. We define the distance from $A$ to $B$ as the length of this path. We denote $d_{T^\rho}(A,B)$ such a distance.
    \item Given $A,B\in \rho$ such that $d_{T^\rho}(A,B)=2$, we let $\overline{e_a}$ and $\overline{e_b}$ be the unique edges of $\overline{T^\rho}$ that correspond to the path from $A$ to $B$ in $\overline{T^\rho}$. We can suppose without loss of generality $a,b$ are both even with $e_a$ and $e_b$ being adjacent to $A$ and $B$ respectively.
    \item Given $A,B\in \rho$, we let $\rho_{A,B}\in \cP(2q)$ be the partition whose blocks are the blocks of $\rho$ except $A$ and $B$ which become the block $A\cup B$ of $\rho_{A,B}$. 

For instance one can see Example on page 6, before Theorem \ref{Lemma: Expansion of moments 5}. 
%For instance, let \(q=3\) and consider
% \[
% \rho
% =
% \bigl\{
% \{1,5\},\{3\},\{2,4\},\{6\}
% \bigr\}
% \in\mathcal P(6).
% \]
% In this case, \(T^\rho\) is a double tree whose underlying tree is
% the path
% \[
% \{3\}
% \;-\;
% \{2,4\}
% \;-\;
% \{1,5\}
% \;-\;
% \{6\}.
% \]
% Let
% \(
% A=\{3\}
% \) and \(
% B=\{1,5\}.
% \)
% Since the unique path from \(A\) to \(B\) is
% \[
% A\;-\;\{2,4\}\;-\;B,
% \]
% we have
% \[
% d_{T^\rho}(A,B)=2.
% \]
% Consequently,
% $$e_a=e_2 \text{ and }e_b=e_4,$$
% and,
% \begin{align*}
% \rho_{A,B}
% &=
% \bigl(\rho\setminus\{A,B\}\bigr)\cup\{A\cup B\}=
% \bigl\{
% \{1,3,5\},\{2,4\},\{6\}
% \bigr\}.
% \end{align*}
% Thus, \(\rho_{A,B}\) is obtained from \(\rho\) by merging \(A\) and
% \(B\) into a single block, while leaving all the other blocks
% unchanged.
\end{enumerate}
\end{notation}

\begin{proposition}\label{Proposition: Infinitesimal contribution pi=q for 2-4 tree}
The following equation satisfies,
\begin{multline*}
\frac{1}{n^{q+1}}\sum_{\substack{\pi\in \cP(2q)\\ T^\pi \text{is a 2-4 tree} }}\sum_{\substack{i_{odd}=1,\dots,n \\ i_{even}=1,\dots,p \\ ker(i)=\pi}}\E(x_i)= 
\sum_{\substack{\sigma\in \NC(2,\dots,2q) \\ \sigma/2\leq ker(r) \\ \rho=\sigma\cup Kr(\sigma)}}\frac{1}{n}\left(\frac{p}{n}\right)^{\#(\sigma)-1}\Bigg(\sum_{\substack{A,B\in \sigma \\ d_{T^\rho}(A,B)=2}}\alpha_{a,b}+\sum_{\substack{A,B\in Kr(\sigma) \\ d_{T^\rho}(A,B)=2}}\frac{p}{n}\alpha_{a,b}\Bigg) + O(\frac{1}{n^2}).
\end{multline*}
Here $\alpha_{a,b}=\frac{1}{2}\E|x_{1,2}|^4 \I_{r_{\frac{a}{2}}=r_{\frac{b}{2}}}+\I_{r_{\frac{a}{2}}\neq r_{\frac{b}{2}}}$, where $a,b$ depend on $A,B$ and are defined as in Notation \ref{Notation: Distance in double tree}.
\end{proposition}

\begin{proof}
Let $\pi\in \cP(2q)$ be such that $T^\pi$ is a 2-4 tree. On the one hand, by Lemma \ref{Lemma: Characterization of 2-4 tree}, $T^\pi$ can be obtained from merging two vertices at distance $2$ of $T^\rho$ where $T^\rho$ is a double tree. Moreover $\rho=\gamma\overline{\rho}$ with $\overline{\rho}=\overline{\pi}^\prime \in \NC_2(2q)$. Conversely, merging two vertices at distance $2$ from a double tree produces a 2-4 tree. 
On the other hand, if $i_1,\dots,i_{2q}$ is such that $ker(i)=\pi$ and $\E(x_i)\neq 0$. From Proposition \ref{Proposition: contribution of 2-4 tree} we know either $\overline{\pi}^\prime \leq ker(rr)$ and $\E(x_i)=1$ or $\overline{\pi}\leq ker(rr)$ and $\E(x_i)=\E|x_{1,2}|^4$. In the first case we get $\overline{\rho}=\overline{\pi}^\prime \leq ker(rr)$, further, there is only one $\overline{\pi}^\prime$ that satisfies this. In the second case we get $\overline{\rho}=\overline{\pi}^\prime \leq ker(rr)$ for the two choices of $\overline{\pi}^\prime$. Thus we can rewrite the expression
$$\sum_{\substack{\pi\in \cP(2q)\\ T^\pi \text{is a 2-4 tree} }}\sum_{\substack{i_{odd}=1,\dots,n \\ i_{even}=1,\dots,p \\ ker(i)=\pi}}\E(x_i),$$
as a sum indexed by $\overline{\rho}\in \NC_2(2q)$ with $\overline{\rho}\leq ker(rr)$, $\rho=\gamma\overline{\rho}$ and $\pi=\rho_{A,B}$ for two vertices $A,B$ at distance two in $T^\rho$. We just need to consider two cases. If $e_a,e_b$ are such that $r_{\frac{a}{2}}=r_{\frac{b}{2}}$, then merging the vertices produces a 2-4 tree where $\overline{\pi}\leq ker(rr)$ and $\E(x_i)=\E|x_{1,2}|^4$. In this case there is another $\overline{\rho} \leq ker(rr)$ that also produces the same $\pi$ and hence we consider both pairings with a contribution of $\frac{1}{2}\E|x_{1,2}|^4$ each. If $e_a,e_b$ are such that $r_{\frac{a}{2}} \neq r_{\frac{b}{2}}$ then merging the vertices produces a 2-4 tree where $\overline{\pi}^\prime \leq ker(rr)$ and $\E(x_i)=1$. In this case there is only one $\overline{\rho} \leq ker(rr)$. Hence,
\begin{equation*}
\sum_{\substack{\pi\in \cP(2q)\\ T^\pi \text{is a 2-4 tree} }}\sum_{\substack{i_{odd}=1,\dots,n \\ i_{even}=1,\dots,p \\ ker(i)=\pi}}\E(x_i)= \sum_{\substack{\overline{\rho}\in \NC_2(2q) \\ \rho=\gamma\overline{\rho} \\ \overline{\rho}\leq ker(rr)}}\sum_{\substack{A,B\in\rho \\ d_{T^\rho}(A,B)=2}}\sum_{\substack{i_{odd}=1,\dots,n \\ i_{even}=1,\dots,p \\ ker(i)=\rho_{A,B}}}\alpha_{a,b}.
\end{equation*}
We now appeal to the same arguments used in Theorem \ref{Lemma: Expansion of moments 4}. $\rho$ can be written as $\rho=\sigma\cup Kr(\sigma)$ with $\sigma\in \NC(\{2,\dots,2q\})$, further from Proposition \ref{Proposition: Equivalence of ker(r) and Ker(rr)} the condition $\overline{\rho}\leq ker(rr)$ becomes $\sigma/2 \leq ker(r)$. Thus
\begin{equation}\label{Aux 8}
\sum_{\substack{\pi\in \cP(2q)\\ T^\pi \text{is a 2-4 tree} }}\sum_{\substack{i_{odd}=1,\dots,n \\ i_{even}=1,\dots,p \\ ker(i)=\pi}}\E(x_i)= \sum_{\substack{\sigma\in \NC(\{2,\dots,2q\}) \\ \rho=\sigma\cup Kr(\sigma) \\ \sigma/2\leq ker(r)}}\sum_{\substack{A,B\in\rho \\ d_{T^\rho}(A,B)=2}}\sum_{\substack{i_{odd}=1,\dots,n \\ i_{even}=1,\dots,p \\ ker(i)=\rho_{A,B}}}\alpha_{a,b}.
\end{equation}
Note that if $A,B\in \rho$ are such that $d_{T^\rho}(A,B)=2$ then either $A,B\in \sigma$ or $A,B\in Kr(\sigma)$ because edges of $T$ always go from even to odd numbers or vice versa. Hence we can simplify the inner sum in the right hand side of Equation \ref{Aux 8} as follows
\begin{eqnarray*}
&&\sum_{\substack{A,B\in \rho \\ d_{T^\rho}(A,B)=2}}\sum_{\substack{i_{odd}=1,\dots,n \\ i_{even}=1,\dots,p \\ ker(i)=\rho_{A,B}}}\alpha_{a,b} \\
&=& \sum_{\substack{A,B\in \sigma \\ d_{T^\rho}(A,B)=2}}\sum_{\substack{i_{odd}=1,\dots,n \\ i_{even}=1,\dots,p \\ ker(i)=\rho_{A,B}}}\alpha_{a,b} + \sum_{\substack{A,B\in Kr(\sigma) \\ d_{T^\rho}(A,B)=2}}\sum_{\substack{i_{odd}=1,\dots,n \\ i_{even}=1,\dots,p \\ ker(i)=\rho_{A,B}}}\alpha_{a,b} \\
&=& \sum_{\substack{A,B\in \sigma \\ d_{T^\rho}(A,B)=2}}\alpha_{a,b}p^{\#(\sigma)-1}n^{\#(Kr(\sigma))} + \sum_{\substack{A,B\in Kr(\sigma) \\ d_{T^\rho}(A,B)=2}}\alpha_{a,b}p^{\#(\sigma)}n^{\#(Kr(\sigma))-1}+O(n^{q-1}), 
\end{eqnarray*}
where the last equality follow from
\begin{eqnarray*}
\sum_{\substack{i_{odd}=1,\dots,n \\ i_{even}=1,\dots,p \\ ker(i)=\rho_{A,B}}}1 &=& (p)_{\#(\sigma)-1}(n)_{\#(Kr(\sigma))}+O(n^{\#(\sigma)+\#(Kr(\sigma))-2}) \\
&=& p^{\#(\sigma)-1}n^{\#(Kr(\sigma))}+O(n^{q-1})
\end{eqnarray*}
whenever $A,B\in \sigma$. A similar statement holds when $A,B\in Kr(\sigma)$. Substituting this into Equation \ref{Aux 8} yields
\begin{multline}
\sum_{\substack{\pi\in \cP(2q)\\ T^\pi \text{is a 2-4 tree} }}\sum_{\substack{i_{odd}=1,\dots,n \\ i_{even}=1,\dots,p \\ ker(i)=\pi}}\E(x_i)= \\
\sum_{\substack{\sigma\in \NC(\{2,\dots,2q\}) \\ \sigma/2\leq ker(r) \\ \rho=\sigma\cup Kr(\sigma)}}\Bigg(\sum_{\substack{A,B\in \sigma \\ d_{T^\rho}(A,B)=2}}\alpha_{a,b}p^{\#(\sigma)-1}n^{\#(Kr(\sigma))} + \sum_{\substack{A,B\in Kr(\sigma) \\ d_{T^\rho}(A,B)=2}}\alpha_{a,b}p^{\#(\sigma)}n^{\#(Kr(\sigma))-1}+O(n^{q-1})\Bigg)
% = \sum_{\substack{\sigma\in \NC(\{2,\dots,2q\}) \\ \sigma/2\leq ker(r) \\ \rho=\sigma\cup Kr(\sigma)}}\left(\sum_{\substack{A,B\in \sigma \\ d_{T^\rho}(A,B)=2}}\alpha_{a,b}\frac{1}{n}\left(\frac{p}{n}\right)^{\#(\sigma)-1} + \sum_{\substack{A,B\in Kr(\sigma) \\ d_{T^\rho}(A,B)=2}}\alpha_{a,b}\frac{1}{n}\left(\frac{p}{n}\right)^{\#(\sigma)}\right)+O(\frac{1}{n^2}),
\end{multline}
Thus, we obtain
\begin{align*}
&\frac{1}{n^{q+1}}\sum_{\substack{\pi\in \cP(2q)\\ T^\pi \text{is a 2-4 tree} }}\sum_{\substack{i_{odd}=1,\dots,n \\ i_{even}=1,\dots,p \\ ker(i)=\pi}}\E(x_i) \\ 
&= \sum_{\substack{\sigma\in \NC(\{2,\dots,2q\}) \\ \sigma/2\leq ker(r) \\ \rho=\sigma\cup Kr(\sigma)}}\Bigg(\sum_{\substack{A,B\in \sigma \\ d_{T^\rho}(A,B)=2}}\alpha_{a,b}\frac{1}{n}\left(\frac{p}{n}\right)^{\#(\sigma)-1} + \sum_{\substack{A,B\in Kr(\sigma) \\ d_{T^\rho}(A,B)=2}}\alpha_{a,b}\frac{1}{n}\left(\frac{p}{n}\right)^{\#(\sigma)}\Bigg)+O(\frac{1}{n^2}),
\end{align*}
which simplifies to the desired result.
\end{proof}

\end{document}
